\section{Contribution comparison, provenance and review navigation}\label{app:comparison}
\subsection{Established mechanisms and the claims submitted here}
Bernstein arithmetic itself is established. Farouki and Rajan \cite{FaroukiRajan} develop arithmetic and other polynomial operations directly in Bernstein form. Their primary publication record also cautions against interpreting basis changes or degree elevation as a general numerical-stability guarantee. The positive convolution in Theorem~\ref{thm:filter} is used as that existing algebra, not claimed as a new arithmetic operation. Its role here is to retain the entire controlled likelihood, including the failed factors and the correct apparatus normalization.

Petrone and Wasserman \cite{PetroneWasserman} study Bernstein polynomial posteriors in Bayesian nonparametric inference and the established representation of Bernstein densities by beta mixtures. Our finite-budget static-parameter filter is a different problem; the beta-mixture identity alone is not the novelty claim. Theorems~\ref{thm:rank} and \ref{thm:amplification} instead identify which physical coefficient functions make the attainable family full dimensional, quantify its interior and degeneracy, and compare passive with selectable likelihood geometry.

The moment-body support calculation belongs to vector-integral convex geometry, for which the lift-zonoid work of Kulik and Tymoshkevych \cite{KulikTymoshkevych} is relevant background. It is not a new convex support principle. Feinberg, Kasyanov and Zgurovsky \cite{FeinbergEtAl} give posterior-state continuity and optimal-policy results for partially observed expected-total-cost problems. The multiplicative payoff here is not that criterion; our compactness and zero-evidence argument verifies the needed properties directly rather than treating their theorem as a plug-in proof.

The endpoint theorem and the linear finite-budget model-error comparison were explicitly proposed by the v3 referee \cite{RefereeV3}. They are credited there and here. The additional control result is Theorem~\ref{thm:arc}: a solved positive-cost gate, a closed circular-arc formula, and four-candidate final-stage mode optimization. Proposition~\ref{prop:example} certifies a strict gap against every report-blind gate, including a full-support prior perturbation. This is a one-stage state-dependent control advantage, not a claimed temporal adaptivity gap over all multi-stage nonadaptive designs.

Theorem~\ref{thm:hierarchy} combines a positive approximation operator, its binomial error estimate, physical detector realization, exact inference and complete-policy regret. The point is not a new Bernstein approximation inequality in isolation: it is an explicit hierarchy that preserves the counted mechanical flags and legal common gates while assigning degree and finite lookup requirements to a desired policy loss. Together with the rank and solved-control theorems, this changes the principal mathematical object from one engineered cubic example to a classified and quantitatively approximable experiment family.

These comparisons use the cited primary publication records and abstracts and the complete internal report; they are not an exhaustive external priority search or a proof audit of those external papers. The manuscript submits the displayed structural claims for independent assessment. Neither a theorem count nor executed diagnostics establishes an editorial judgment about mathematical significance.

\subsection{Pinned controlling sources and actual historical reuse}
The controlling English-v3 review is on
\begin{quote}\small\path{review/a1-english-v3-harsh-referee-2026-09-05}\end{quote}
at commit \path{574f2315a136d8b401644d8a3eeeb93c87887010}. The reviewed manuscript commit and its source tree are, respectively:
\begin{quote}\small
\path{025a9b3fdfd9ffa86e6ae2924f8b5e5b1b58ddd6}\\
\path{ec09fa385c5fdcc5f2aece59a6b7f0806bfbb8ea}.
\end{quote} The report's two local formulation objections concern smooth reporting maps and intrinsic moment-body dimension; they are addressed in Theorem~\ref{thm:chamber-repair} and Remark~\ref{rem:dimension}. The established corrections to biased entropy, two physical interface traces and explicit preparation transport remain in Appendix~\ref{app:repairs}.

The complete principal v3 source, referee report and claim audit were read. Focused historical reads included its companion Sections 13 and 16, on hybrid sensitivity and the whole-preparation collision-tube identity, the repository's \path{ROUND12_HISTORICAL_DERIVATION_AUDIT.md}, and the local primitive/pointer part of \path{papers/A1-exact-benchmarks/ROUND17_POSITIVE_CLOSURE.tex}. The latter primitive is checked directly: for
\[
 (Q,P)=((q-s)/w,s+wp),\qquad P\,dQ-p\,dq=d\{s(q-s)/w\}.
\]
This and the elementary pointer impulse motivate only the local constructions used here. No historical global quotient, conormal calculus, spectral, renewal or central-limit claim is imported as an unverified input.

The historical whole-preparation argument supplies a complete collision-tube change of variables and separates physical radius-dependent preparation from fixed ambient preparation. The present counted-cartridge normalization is rederived in Section~\ref{sec:framework}. The new rank, gate and approximation theorems do not assume an unproved many-collision hierarchy. The companion's earlier general material remains available at its pinned source; only the specifically cited parts are represented as freshly examined, not all 63 historical statements or all eleven papers.

\subsection{Reproduction and the next referee object}
This revision is an ordinary, self-contained LaTeX source tree. Its adjacent response letter maps M1, M2, E1--E3, G1 and S1 to theorem labels and describes the retained earlier repairs. The proof ledger records dependencies and exact quantifiers. The test directory contains finite algebraic/numerical regression diagnostics and a separate rational interval certificate for Proposition~\ref{prop:example}; validation receipts record actual execution and compilation, with hashes.

Finite rank checks do not replace the all-$n$ product-differential proof, finite gates do not verify compact infinite-dimensional selection, and numerical sampling does not certify collision-tube coverage. The rational interval calculation does certify its stated numerical inequalities because it uses outward rational enclosures and explicit transcendental remainder bounds. Independent review should assess the theorem proofs and their significance separately from these reproducibility checks. Existing main, review and previous revision sources are preserved.
